\section{Mixed volumes}
\label{sec-mixed}

Let $H$ be a standard graded, Artinian, Gorenstein $K$-algebra of socle degree $n$,
\[ H = K[x_1,\ldots,x_N]/I. \]
It is well-known that $H$ is determined by the linear function
\[  W: K[x_1,\ldots,x_N]_n \to H^n \stackrel{\isom}{\longrightarrow} K.\]
(We have denoted by subscript $n$ the degree $n$ homogeneous part of $K[x_1,\ldots,x_N]$.) Indeed,  one recovers the ideal $I$ from $W$ using the property that $f\in K[x_1,\ldots,x_N]_m$ lies in $I$ if and only if $W(fg) = 0$ for any $g$ of degree $n-m$. More generally, any nonzero linear function $W: K[x_1,\ldots,x_N]_n \to K$ determines a standard graded, Artinian, Gorenstein $K$-algebra $H$ of socle degree $n$.

For an oriented pseudo-manifold $\Delta$, let the function $W = W_\Delta$ be the composition
\[ W_\Delta:  K[x_1,\ldots,x_N]_n \to \cA^n(\Delta) \stackrel{\pi_\Delta}{\longrightarrow} K.\]
When $\Delta$ is a homology sphere over $K$, then $W_\Delta$ determines the algebra $H(\Delta)$. When $\Delta$ is an oriented pseudo-manifold over $K$, then $W_\Delta$ determines an algebra that we denote $\bar{H}(\Delta)$. This algebra in general is a quotient of the algebra $H(\Delta)$.

In the theory of polytopes and toric varieties the function $W_\Delta$ is known as the mixed volume.


\subsection{The case of $\Pi$} \label{sec-Pi}
Let $\pi$ be an $n$-simplex and $\Pi = \partial\pi$ the $(n-1)$-dimensional sphere. We compute here the mixed volume $W_\Pi$. 

Let $v_0, v_1,\ldots, v_n$ be the vertices of $\Pi$, and $\sigma_j = \{v_0,\ldots, \hat{v}_j, \ldots, v_n\}$ the maximal simplices. We choose the orientation on $\Pi$ so that $v_1,\ldots,v_n$ is positively oriented on the simplex $\sigma_0$.   Denote by 
\[ A = (\a{i, j})_{i,j} \]
the $n\times (n+1)$ matrix of variables, where the columns are indexed by $0,1,\ldots, n$ and the rows by $1,\ldots,n$. Let $X_j \in K$ be $(-1)^j$ times the determinant of the matrix $A$ with its $j$-th column removed. Then $X_j = \det \sigma_j$ as defined in Equation~\eqref{eq-dets} on page~\pageref{eq-dets}.

\begin{lemma} \label{lem-Pi}
Let $f\in K[x_0,\ldots, x_n]_n$. Then
\[ W_\Pi \big( f(x_0,x_1,\ldots,x_n) \big) =   \frac{f(X_0, X_1, X_2, X_3,\ldots, X_n)}{X_0 X_1\cdots X_n}.\]
\end{lemma}

\begin{proof}
We first check that $W_\Pi(\theta_i g) = 0$ for any $g$ of degree $n-1$ and $i=1,\ldots,n$. It suffices to show that $\theta_i$ evaluated at $X_0,\ldots, X_n$ is zero. From the definition,
\[ \theta_i(X_0,\ldots, X_n) = \sum_{j=0}^n  \a{i, j} X_j.\]
This sum is the expansion of the determinant of the matrix $A$ with a copy of its $i$-th row added as the first row. Since the matrix has two repeated rows, its determinant is zero. 

The previous argument shows that the map $W_\Pi$ factors through $H^n(\Delta)$.  Let us check that its value on the monomial $\chi_{\sigma_0} = (\det\sigma_0) x_1\cdots x_n$ is $1$ as required:
\[ W_\Pi((\det\sigma_0)x_1\cdots x_n) = \frac{X_0 X_1 \cdots X_n}{X_0 X_1\cdots X_n} = 1. \qedhere\]
\end{proof}

To simplify notation, let us write the mixed volume as
\[ W_\Pi = \frac{1}{c_\Pi} \ev_\Pi,\]
where $c_\Pi = X_0 X_1 \cdots X_n \in K$ and $\ev_\Pi: K[x_0,\ldots, x_n] \to K$ is the evaluation map that sets $x_j= X_j$. The evaluation map is a $K$-algebra homomorphism.

Recall that in Section~\ref{sec-connected-sum} we decomposed an oriented pseudo-manifold $\Delta$ as a connected sum
\[ \Delta = \#_{i=1}^M \Pi_i.\]
The following result now follows from Lemma~\ref{lem-int-sum} and Lemma~\ref{lem-Pi}:

\begin{theorem} \label{thm-W-sum}
Let $\Delta$ be an oriented pseudo-manifold. Then
\[ W_\Delta = \sum_{i=1}^M W_{\Pi_i} = \sum_{i=1}^M \frac{1}{c_{\Pi_i}} \ev_{\Pi_i}.\]
\end{theorem}

In the theorem the map $\ev_{\Pi_i}$ acts on $K[x_1,\ldots,x_N]$ as a composition
\[ K[x_1,\ldots, x_N] \to K[x_{j_1},\ldots, x_{j_n}] \stackrel{\ev_{\Pi_i}}{\longrightarrow} K,\]
where the first map, the restriction to $\Pi_i$, sets $x_j=0$ if $v_j$ does not lie in $\Pi_i$. When $v_j = v_{j_l}$ is a vertex of $\Pi_i$, then $\ev_{\Pi_i}$ maps $x_j$ to $X_j$. However, the constants $X_j\in K$ depend not only on $j$ but also on all vertices of $\Pi_i$ and the orientation on $\Pi_i$.
